\chapter{Case Study: Direct N-Body Simulation}
\label{ch:cuda-nbody}

The previous case studies focused on computations with dependencies between
updates. An $N$-body simulation provides a different form of parallelism: at a
given time step, the acceleration of each body can be computed independently
from the current positions of all bodies. This gives a simple, useful example
of an all-pairs computation on the GPU.

This chapter starts with direct summation, in which every body interacts with
every other body. It then improves the force calculation by loading source
bodies into shared memory and reusing them across the threads of a block. The
number of interactions remains $O(N^2)$; the optimization targets data reuse.
The companion notebook, \lstinline|chapter07_nbody_direct_cuda.ipynb|,
contains the complete CUDA program, a CPU reference, timing code, and a Python
animation of the simulation.

\begin{learningbox}
The reader will learn:
\begin{itemize}
    \item how to map an all-pairs force calculation to one CUDA thread per
          target body;
    \item how to reuse source-body data with a shared-memory tile;
    \item how to organize a Velocity-Verlet step as ordered CUDA kernels;
    \item how to validate GPU results against a CPU reference;
    \item how to measure force-kernel time separately from the complete
          simulation and visualize the resulting motion.
\end{itemize}
\end{learningbox}

\section{The Gravitational Model}
\label{sec:nbody-model}

Let body $i$ have position $\mathbf{x}_i$, velocity $\mathbf{v}_i$, and mass
$m_i$. The acceleration of body $i$ due to all the other bodies is
\begin{equation}
\mathbf{a}_i = G \sum_{\substack{j=0\\j\ne i}}^{N-1} m_j
\frac{\mathbf{x}_j-\mathbf{x}_i}
     {\left(\|\mathbf{x}_j-\mathbf{x}_i\|^2+\varepsilon^2\right)^{3/2}}.
\label{eq:nbody-force}
\end{equation}
The softening length $\varepsilon$ prevents the force from becoming singular
when two bodies are very close. The notebook uses normalized units with
$G=1$ and $\varepsilon=0.02$. Its initial configuration consists of a massive
body at the center and a much lighter disk of bodies in the $xy$ plane. The
force calculation itself is three-dimensional. Disk velocities are chosen to
be approximately circular in the field of the central body; this is a
teaching model, rather than a calibrated astrophysical system.

The notebook advances the state with the Velocity-Verlet method. Given the
position, velocity, and acceleration at time step $n$, first compute
\begin{equation}
\mathbf{x}_i^{n+1} = \mathbf{x}_i^n + \Delta t\,\mathbf{v}_i^n
             + \frac{1}{2}\Delta t^2\mathbf{a}_i^n.
\label{eq:nbody-position-update}
\end{equation}
The accelerations are then evaluated at the new positions, and the velocities
are updated as
\begin{equation}
\mathbf{v}_i^{n+1} = \mathbf{v}_i^n +
   \frac{\Delta t}{2}\left(\mathbf{a}_i^n+\mathbf{a}_i^{n+1}\right).
\label{eq:nbody-velocity-update}
\end{equation}
This method requires two force evaluations for the first step (one to obtain
the initial accelerations and one at the new positions) and one new force
evaluation per subsequent step. The notebook initializes the first
acceleration on the CPU, then keeps the evolving state on the GPU.

For each target body, the force loop visits the other $N-1$ bodies. The total
work per time step is therefore proportional to $N(N-1)$, or $O(N^2)$. The
force on each target is an independent output, so there are no writes shared
between target threads and no atomics are needed.

\section{Body Data and CPU Reference}
\label{sec:nbody-data}

The CUDA program stores each body's position and mass together in a
\lstinline|double4_16a| record: the \lstinline|x|, \lstinline|y|, and
\lstinline|z| fields hold position, and \lstinline|w| holds mass. Velocity
and acceleration are stored in separate arrays of the same record type; their
fourth fields are unused. Neighboring CUDA threads process neighboring body
indices, so their accesses to these arrays are adjacent. The explicit
16-byte-aligned vector type avoids the deprecated \lstinline|double4| name in
CUDA Toolkit 13 and later.

A CPU implementation of the same force law and integrator provides a
reference for validation. The GPU and CPU sum interactions in different
orders, so their floating-point results need not be bitwise identical. The
notebook compares the state after one step using a maximum absolute error and
a tolerance of $10^{-9}$.

\section{Direct Force Evaluation}
\label{sec:nbody-naive}

The baseline GPU kernel assigns one thread to each target body. The thread
loads that body's position, loops over all source bodies, accumulates the
three acceleration components, and writes one acceleration record. The
integrator is kept separate from this force kernel so the same update kernels
can be used with both force implementations. The CUDA source uses $G=1$ and
a block size of 256:

\begin{lstlisting}[language=C++]
constexpr int BLOCK_SIZE = 256;
constexpr double GRAVITATIONAL_CONSTANT = 1.0;
\end{lstlisting}

\begin{lstlisting}[language=C++,caption={Direct all-pairs acceleration: one thread computes the acceleration of one target body.},label={lst:nbody-naive}]
__global__ void acceleration_global(
    const double4_16a* __restrict__ position_mass,
    double4_16a* __restrict__ acceleration,
    int n,
    double softening_squared)
{
    const int i = blockIdx.x * blockDim.x + threadIdx.x;
    if (i >= n) return;

    const double4_16a target = position_mass[i];
    double ax = 0.0, ay = 0.0, az = 0.0;

    for (int j = 0; j < n; ++j) {
        if (i == j) continue;
        const double4_16a source = position_mass[j];
        const double dx = source.x - target.x;
        const double dy = source.y - target.y;
        const double dz = source.z - target.z;
        const double r2 = dx*dx + dy*dy + dz*dz + softening_squared;
        const double inv_r = 1.0 / sqrt(r2);
        const double scale = GRAVITATIONAL_CONSTANT * source.w * inv_r * inv_r * inv_r;
        ax += scale * dx;
        ay += scale * dy;
        az += scale * dz;
    }
    acceleration[i] = make_double4_16a(ax, ay, az, 0.0);
}
\end{lstlisting}

The bounds check handles the extra threads in the final block when $N$ is not
a multiple of the block size. Each valid thread writes a distinct output
index. At a given loop iteration, the lanes of a warp read the same source
index $j$; the hardware caches may therefore capture some of the reuse even
in this direct version. The tiled kernel makes reuse across a whole block
explicit.

\section{Reusing Source Bodies with Shared Memory}
\label{sec:nbody-tiled}

The tiled kernel also assigns one thread to each target body. At each loop
iteration, however, the block cooperatively loads a contiguous tile of source
records into shared memory. Every valid target thread uses that tile to
accumulate part of its acceleration before the block loads the next tile.
The notebook uses a block size of 256 and a statically sized shared-memory
array.

\begin{lstlisting}[language=C++,caption={Tiled all-pairs acceleration with shared-memory reuse.},label={lst:nbody-tiled}]
__global__ void acceleration_tiled(
    const double4_16a* __restrict__ position_mass,
    double4_16a* __restrict__ acceleration,
    int n,
    double softening_squared)
{
    __shared__ double4_16a source_tile[BLOCK_SIZE];
    const int i = blockIdx.x * blockDim.x + threadIdx.x;
    const int lane = threadIdx.x;
    double ax = 0.0, ay = 0.0, az = 0.0;
    double4_16a target = make_double4_16a(0.0, 0.0, 0.0, 0.0);
    if (i < n) target = position_mass[i];

    const int tile_count_total = (n + BLOCK_SIZE - 1) / BLOCK_SIZE;
    for (int tile = 0; tile < tile_count_total; ++tile) {
        const int source_index = tile * BLOCK_SIZE + lane;
        if (source_index < n)
            source_tile[lane] = position_mass[source_index];
        else
            source_tile[lane] = make_double4_16a(0.0, 0.0, 0.0, 0.0);
        __syncthreads();

        const int count = min(BLOCK_SIZE, n - tile * BLOCK_SIZE);
        if (i < n) {
            for (int k = 0; k < count; ++k) {
                const int j = tile * BLOCK_SIZE + k;
                if (i == j) continue;
                const double4_16a source = source_tile[k];
                const double dx = source.x - target.x;
                const double dy = source.y - target.y;
                const double dz = source.z - target.z;
                const double r2 = dx*dx + dy*dy + dz*dz
                                  + softening_squared;
                const double inv_r = 1.0 / sqrt(r2);
                const double scale = GRAVITATIONAL_CONSTANT * source.w * inv_r * inv_r * inv_r;
                ax += scale * dx;
                ay += scale * dy;
                az += scale * dz;
            }
        }
        __syncthreads();
    }
    if (i < n)
        acceleration[i] = make_double4_16a(ax, ay, az, 0.0);
}
\end{lstlisting}

Both barriers must be reached by every thread in the block. The first makes
the loaded tile available to the block; the second ensures that all target
threads have finished reading it before any thread overwrites it with the
next tile. Threads in a partially filled final block skip the force
accumulation when their target index is invalid, but still participate in
both barriers.

A source tile of $B$ records is loaded once per block and reused by up to $B$
target threads. This reduces repeated source loads at the block level, while
the number of target--source interactions remains $O(N^2)$. The reduction in
source-level loads does not translate directly into the same reduction in
DRAM traffic: the direct kernel's same-index warp loads may be served
efficiently by caches or broadcast behavior. The notebook therefore measures
both versions on the available GPU rather than assuming that shared memory
always wins.

\section{A Velocity-Verlet Step on the GPU}
\label{sec:nbody-device-loop}

One time step is implemented as three ordered kernels:
\begin{enumerate}
    \item update all positions using the current velocities and accelerations;
    \item compute new accelerations from the updated positions, using either
          \lstinline|acceleration_global| or
          \lstinline|acceleration_tiled|;
    \item update all velocities and replace the stored accelerations by the
          new values.
\end{enumerate}
The kernels are launched in the same CUDA stream. Their order provides the
global ordering between these stages: the force kernel sees all updated
positions, and the velocity kernel sees all new accelerations. No host-side
synchronization is needed between steps. The evolving positions, velocities,
and accelerations stay in device memory; the host copies positions back only
when saving frames for the animation.

\section{Validation, Timing, and Visualization}
\label{sec:nbody-validation}

The notebook first initializes a small system with $N=64$, computes one
Velocity-Verlet step on the CPU, and compares it with one step using each GPU
force kernel. The maximum absolute error is checked against $10^{-9}$. This
checks both force implementations and the common integration stages.

For timing, the notebook uses $N=2048$, performs three warm-up launches, and
averages ten launches measured with CUDA events. Allocation and host-to-device
transfer occur outside the timed interval, so the reported value is force
kernel time per launch, not a full time step or an end-to-end simulation
measurement. Results depend on the GPU, CUDA version, and runtime conditions;
the notebook prints the measurements obtained on the reader's device.

The animation uses the tiled kernel for $N=512$ bodies and 1200 time steps,
with $\Delta t=0.001$. It saves the initial state and then one snapshot every
10 steps, giving 121 frames. A short Python cell reads the positions and
animates the disk and central body. Python is used only for visualization;
the force calculation and time integration are written in CUDA C++.

\begin{validationbox}[Common pitfalls]
\begin{itemize}
    \item Returning from a tiled kernel before a block-wide barrier. All
          threads must reach both \lstinline|__syncthreads()| calls.
    \item Omitting the second barrier, which may let a thread overwrite a
          source tile while another thread is still using it.
    \item Updating each body's position in the same kernel that computes
          forces from positions that other threads are concurrently reading.
          The notebook separates the position update and force evaluation
          into ordered launches.
    \item Setting the softening parameter to zero, which permits a singular
          force for coincident bodies.
    \item Treating one kernel's timing as the run time of the complete
          multi-step simulation.
\end{itemize}
\end{validationbox}

\section{Summary and Further Directions}
\label{sec:nbody-summary}

Direct summation exposes $N$ independent target calculations, each interacting
with all $N$ sources. The baseline kernel assigns one thread to each target;
the tiled kernel reuses source records in shared memory. Both perform
$O(N^2)$ interactions. The Velocity-Verlet update is split into ordered GPU
kernels, and the accompanying notebook validates the calculation, measures
the force kernels, and animates the resulting motion.

For much larger systems, treecodes and the fast multipole method (FMM) reduce
the work by approximating interactions between sufficiently distant groups of
bodies. They require additional data structures and algorithms, so they are
outside the scope of this direct-method case study and may be treated as a
later chapter or supplementary example.

\begin{enumerate}
\item For $N=1000$ and a block size $B=256$, how many target threads are
      launched? How many source tiles are visited by each valid target?
\item Why must invalid target threads in the final block still reach both
      barriers in \lstinline|acceleration_tiled|?
\item Shared-memory tiling reduces repeated source loads but does not change
      the number of interactions. Explain why these statements are
      compatible.
\item Which parts of the program are included in the notebook's force-kernel
      timing? What additional work would need to be measured for an
      end-to-end simulation time?
\end{enumerate}

\section*{References}
\addcontentsline{toc}{section}{References}
\begin{enumerate}
\item NVIDIA, \emph{CUDA C++ Programming Guide}, sections on vector types,
      shared memory, and synchronization.
\item NVIDIA, \emph{CUDA C++ Best Practices Guide}, sections on coalescing,
      shared-memory reuse, and performance measurement.
\end{enumerate}
